\section{Introduction}
\label{sec:intro}
Evaluating a language model on a modern benchmark means answering thousands of items; repeating it for many checkpoints, prompts or inference settings is expensive. A tempting shortcut is to run a model on a small subset and extrapolate to the full benchmark, using item-level results of models that have already been evaluated (``source models'') to choose the subset and to calibrate the extrapolation. Several recent methods do this with clustering, item response theory (IRT) or learned selection policies~\citep{vivek2024anchor,maiapolo2024tiny,kipnis2025metabench,li2025active,pacchiardi2024100}.

What should we expect from such methods in the situation where they are most needed, i.e.\ when the model under test is \emph{new} and different from the source models? Two issues make this setting hard. First, a method is typically validated on a random split of models, where target models resemble sources. Second, any dependence on the source models enters through two channels---\emph{which items are kept} and \emph{how the subset is turned into a score}---and these are rarely separated, so it is unclear which channel brings the gains or the failures.

We study this with a controlled factorial design on 13 binary item-level score matrices (Open LLM Leaderboard v2 and HELM, up to 429 models and 12{,}032 items). We cross six item selectors with several estimators, evaluate under four protocols (IID model split, two \emph{extrapolation} splits in which targets are the more accurate models, and a held-out-family split), report MAE, RMSE, Spearman and Kendall with bootstrap intervals over benchmarks and replicates, and measure bias and interval coverage with repeated subsampling.
Because plain model-based estimators turn out to be strongly biased under extrapolation, we additionally study a per-target recalibration of the predictor, a cross-fitted, finite-population difference estimator (CF) that is exactly unbiased under random item sampling, and an empirical-Bayes variant (CF-EB) that shrinks towards the plug-in when the cohort's own data show the plug-in is accurate.

\textbf{Contributions.} (1) A factorial, paired comparison of item selectors and estimators with uncertainty quantification, separating the effect of selection from that of estimation and the IID from the extrapolation regime (\cref{sec:results}). (2) A cross-fitted difference estimator with a proof of exact unbiasedness under random item sampling, design-based confidence intervals, and a cohort-adaptive shrinkage version (\cref{sec:method}). (3) Evidence on when each family works and a list of failure modes, including negative results (\cref{sec:results,sec:discussion}). All code and the pinned data version are provided for reproduction.
